\section{A diagnostic, and what it implies}\label{sec:diag}

Choosing the number of admitted jobs from the throughput curve, by paging heuristics~\citep{denning1976optimal} or from measured throughput~\citep{heiss1991adaptive,schroeder2006mpl}, is the classical selection of a multiprogramming level. Stack-distance estimators size a KV tier to the working set, for a target hit rate (KVSET~\citep{li2026kvset}) or for the host offload tier (EfficientAgent~\citep{shao2026efficientagent}), and Yao et al.'s mean-field estimator sizes memory for a target hit ratio when turn gaps are exogenous~\citep{yao2026multiturn}. What the window adds is that the from-empty edge, a finite-horizon stability limit in the sense of Kleinrock and Lam's first exit time~\citep{kleinrock1975aloha}, and not only $\mu^*$, bounds the operating range; that a cap protects against history, not only throughput; and that the right cap moves with the tool-time law, the token profile and the pool.

The window of Eq.~\eqref{eq:window} comes from three measurements that an operator can make before deployment, on the engine, model, pool and workload that will be served. (1) \emph{All-miss capacity}: run a saturated closed population with prefix caching disabled and measure the uncached prefill token rate, which is stable across seeds; check saturation by measuring at two populations (for example 32 and 64 sessions) and use the larger rate, since a population that is not saturated understates $\lmiss$ (our 32-session reference fell about 2--3\% short, \S\ref{sec:props}); convert it to sessions by the mean work per session, or use turns per second over $E[T]$ measured over a window longer than a session's lifetime at that population (mini-SWE-agent's sessions outlived the declared 1800\,s window, and the turn rate there gave a $\lmiss$ 15\% low); completed sessions over a short window are noisy. By Proposition~\ref{prop:A} and the fluid model of \S\ref{sec:fluid}, under request-level FCFS this is an upper bound on the recovery threshold from deep congestion if the congested engine computes prompt tokens no faster than the cache-disabled one, and it equals that threshold if, in addition, $\mu(n)\ge\lmiss$ at every population beyond the peak and not only in the limit (with zero hits, if the engine computes prompt tokens as fast as the cache-disabled one at every congested population); the heuristic work balance of \S\ref{sec:fluid} leads to the same conditions for the open system. With zero hits our engine computed prompt tokens at 1.03 times the all-miss rate near the stationary turn mix and at 0.99--1.12 times it over the windows we measured, so (A3) and (A4) hold on it only to within a few percent at the stationary mix and to within 12\% over those windows. If sessions share a prefix (system prompt, tool definitions), measure instead with caching on and each turn's private context fresh, so that only the shared prefix can hit, the reference against which we expect Proposition~\ref{prop:A} to hold for the private suffixes (\S\ref{sec:props}); disabling caching entirely understates the congested rate. (2) \emph{Closed-population curve}: measure $\mu(n)$ on a grid that crosses the cliff, and do not read $\muinf$ off the largest $n$ of a short run with a synchronized start; here $\mu(300)$ was 0.108\sps, 51\% above $\lref$. (3) \emph{From-empty edge}: run open arrivals from an empty engine at increasing rates over the horizon of interest, at least that of the burst test.

If $\mu^*>\lmiss$, then in the fluid model of \S\ref{sec:fluid}, under request-level FCFS and the assumptions of Proposition~\ref{prop:A}, every operating point in $(\lmiss,\min(\mu^*,\ledge))$ is one sufficiently large burst away from the congested state unless admission is bounded; a steady-state load test there passes without testing recovery. On the engine we tested two rates per pool and one burst size: 200 sessions, about 20 times the 8--12 present at 300k and 6 to 13 times the 16--34 present with the native pool, left all 12 registered SGLang BURST runs at zero hits (11 trapped by the rule); in the simulator 25 never trapped and 50 trapped about half the time. Without a cap, Proposition~\ref{prop:A} and the fluid model predict that recovering from deep congestion requires cutting the arrival rate below $\lmiss$, which was 31\% of peak throughput at 300k tokens and 15\% with the native pool; the earlier, uncalibrated simulator agrees (over 150,000\,s, FCFS recovered from congestion at 0.9 times its own $\lmiss$ and not at 1.0), but no real run lowered the rate after a trap. Stock LPM, which on our reading of its source stops re-sorting its queue when more than 128 requests wait and so falls under Proposition~\ref{prop:A} in deep congestion, recovered from the same finite burst at 1.35 $\lref$ in 4 of 6 grid runs, and at 1.65 $\lref$ in 1 of 6 (\S\ref{sec:grid}); the proposition is a limit statement and does not predict recovery from a finite burst for any ordering. Proposition~\ref{prop:cap} and the quasi-static model suggest a session cap at or just below the peak, since a cap past the peak leaves $(\mu(c),\mu^*)$ open (\S\ref{sec:props}); we take the smallest $n$ within a few percent of the peak (20 at 300k, 64 with the native pool), because the curve has one run per $n$, and check it with a burst test. We recommend a cap because it needs no engine change, its effect in the quasi-static model follows from Proposition~\ref{prop:cap} and Eq.~\eqref{eq:fluid}, and it is the one mechanism we tested on the engine (all 12 registered SGLang CAPPED runs recovered), not because it dominates: in the earlier simulator, 20\,s TTL pins with a pinned-first order reached a higher $\mu(300)$ than any static cap (\S\ref{sec:mech}). A cap does not make the burst free: it moves the burst's cost to the client as admission delay (per-run maxima of 853--985\,s at 300k and 247--361\,s with the native pool), instead of turns that take several minutes for the rest of the run (\S\ref{sec:e2}, \S\ref{sec:e2native}). The delay ends only if the gate drains, which in the quasi-static model needs $\lambda<\mu(c)$ for a cap $c$ (Proposition~\ref{prop:cap} and Eq.~\eqref{eq:fluid}). Whether the gate should queue or reject sessions is a deployment choice we did not evaluate.

The cap must be measured again when the workload or configuration changes. In the simulator the peak moved from 20 to 32 or 40 sessions with the TraceLab tool-time laws and to 24 with mini-SWE-agent's (where the cap stayed at 20), and fell to 12 with mini-SWE-agent's mean turn and token counts. On the real engine it moved to 64--80 with the 955,757-token pool, where the registered cap of 64 drained the burst in every run, and with vLLM 0.32.0 on the 300k pool it stayed at 20--24. A cap sized for one setting can fail in another in either direction. Too small, it holds the engine below the offered load: a cap of 20 under the TraceLab idle-gap law gives $\mu(20)=1.26\,\lmiss$, below both test rates (1.35 and 1.65 $\lmiss$), so the gate would not drain. Too large, it sits on the collapsed side of the curve: with mini-SWE-agent's full workload $\mu(20)$ is 1.45 times that workload's declared $\lmiss$, again below both test rates (1.65 and 2.35). Neither mis-sized cap was run. The procedure sizes a static cap offline and gives no online detector of a trap in progress. Here its three stages took 14--16 A100-hours per pool and engine, and a three-arm burst check about 7.5 more per rate and seed.
